\documentclass[12pt]{article}
\usepackage[utf8]{inputenc}
\usepackage[T1]{fontenc}
\usepackage{hyperref}
\title{Kulon Language Resources for K-12 Curriculum Development}
\author{Compiled by Claude Sonnet 5 (claude-sonnet-5) for Andr\'as Kornai}
\date{}
\begin{document}
\maketitle

\section{Introduction}
Kulon (uon) is an extinct Austronesian language of Taiwan and is very poorly documented. The main source is a roughly 60-page study by Tsuchida (1985). Its classification is disputed: Li (2008) ties it to Saisiyat, while Blust (1999) links it more closely to Pazeh. It should not be confused with Kulon-Pazeh (uun), which is the code for Pazeh. Pazeh's last native speaker died in 2010, and the related Kaxabu is still being documented. I did not find an extinction date for Kulon.

\section{Primary Educational Resources}
None found. There are no known primers, textbooks or programmes for Kulon.

\section{Linguistic and Reference Materials}
\begin{itemize}
\item Wikipedia, Kulon language (summary of the sources, including Tsuchida 1985): \url{https://en.wikipedia.org/wiki/Kulon_language}
\item Wikipedia, Formosan languages (context): \url{https://en.wikipedia.org/wiki/Formosan_languages}
\item Endangered Languages Project entry: \url{https://www.endangeredlanguages.com/ru/node/96215}
\end{itemize}

\section{Religious and Community Resources}
None found.

\section{Digital and Audio Resources}
None found for Kulon. Audio recordings exist for the related Kaxabu, but I did not find a source for them.

\section{Curriculum Development Context}
Given the very limited documentation, a Kulon revitalization curriculum would be mostly reconstructive. It could rest on Tsuchida (1985), whose text I did not locate online, and on comparison with Pazeh and Saisiyat. Any such work should involve Taiwanese Indigenous community and academic partners. No existing school programme was found.
\end{document}